\section*{Bloque B. Aplicaciones cuantitativas y optimización}

\noindent\textit{Preguntas 31 a 65. Requieren resolver problemas de maximización de beneficios, derivar ofertas individuales y de industria, calcular equilibrios de corto y largo plazo, determinar excedentes y evaluar los efectos de impuestos, subsidios y controles de precios.}

\q{Si $C(q)=100+10q+q^2$ y $P=30$, la cantidad que maximiza el beneficio es:}
\begin{opts}
  \item 5
  \item 10
  \item 15
  \item 20
\end{opts}

\q{En el problema anterior, el beneficio económico en el óptimo es:}
\begin{opts}
  \item 0
  \item 100
  \item $-100$
  \item 50
\end{opts}

\q{Si $C(q)=50+5q+0.5q^2$ y $P=25$, la cantidad que maximiza el beneficio es:}
\begin{opts}
  \item 10
  \item 15
  \item 20
  \item 25
\end{opts}

\q{En el problema anterior, el beneficio económico es:}
\begin{opts}
  \item 100
  \item 150
  \item 200
  \item 250
\end{opts}

\q{Si $C(q)=100+4q+q^2$ y $P=20$, el beneficio económico en el óptimo es:}
\begin{opts}
  \item $-36$
  \item 36
  \item 0
  \item $-64$
\end{opts}

\q{Con $C(q)=100+4q+q^2$, el precio de cierre de la empresa en el corto plazo es:}
\begin{opts}
  \item 0
  \item 2
  \item 24
  \item 4
\end{opts}

\q{Si $C(q)=q^2+10q+50$ y $P=30$, la cantidad ofrecida por la empresa es:}
\begin{opts}
  \item 5
  \item 10
  \item 15
  \item 20
\end{opts}

\q{Si 100 empresas idénticas tienen oferta individual $q=(P-10)/2$, la oferta de la industria es:}
\begin{opts}
  \item $Q=50P-500$
  \item $Q=50P+500$
  \item $Q=100P-1000$
  \item $Q=100P-500$
\end{opts}

\q{Si la demanda de mercado es $Q=1000-50P$ y la oferta es $Q=50P-500$, el precio de equilibrio es:}
\begin{opts}
  \item 15
  \item 10
  \item 20
  \item 25
\end{opts}

\q{En el equilibrio anterior, la cantidad de mercado es:}
\begin{opts}
  \item 100
  \item 150
  \item 200
  \item 250
\end{opts}

\q{En el mismo equilibrio, la cantidad producida por cada una de las 100 empresas es:}
\begin{opts}
  \item 1
  \item 2.5
  \item 5
  \item 10
\end{opts}

\q{Si $C(q)=q^2+4q+100$, el precio de equilibrio de largo plazo es:}
\begin{opts}
  \item 20
  \item 22
  \item 24
  \item 26
\end{opts}

\q{En el equilibrio de largo plazo anterior, la cantidad producida por cada empresa es:}
\begin{opts}
  \item 10
  \item 5
  \item 20
  \item 24
\end{opts}

\q{Si la demanda de mercado es $Q=500-10P$ y cada empresa produce 10 unidades en el largo plazo, el número de empresas es:}
\begin{opts}
  \item 20
  \item 24
  \item 26
  \item 30
\end{opts}

\q{En el equilibrio de largo plazo con $P=24$ y $Q=260$, el excedente del consumidor es:}
\begin{opts}
  \item 2600
  \item 3380
  \item 5000
  \item 6760
\end{opts}

\q{En el equilibrio de largo plazo competitivo con libre entrada, el beneficio económico de cada empresa es:}
\begin{opts}
  \item 3380
  \item 2600
  \item 0
  \item 260
\end{opts}

\q{Si en una industria de costo constante la demanda aumenta, en el nuevo equilibrio de largo plazo el precio:}
\begin{opts}
  \item vuelve a 24.
  \item sube permanentemente a 30.
  \item baja a 20.
  \item sube a 26 y se queda allí.
\end{opts}

\q{Si la demanda es $Q_d=120-2P$ y la oferta es $Q_s=3P-30$, el precio de equilibrio es:}
\begin{opts}
  \item 20
  \item 30
  \item 40
  \item 50
\end{opts}

\q{En el equilibrio anterior, la cantidad de equilibrio es:}
\begin{opts}
  \item 30
  \item 50
  \item 60
  \item 80
\end{opts}

\q{En el mismo mercado, el excedente del consumidor es:}
\begin{opts}
  \item 300
  \item 600
  \item 1200
  \item 900
\end{opts}

\q{En el mismo mercado, el excedente del productor es:}
\begin{opts}
  \item 300
  \item 900
  \item 600
  \item 1200
\end{opts}

\q{En el mismo mercado, el excedente total es:}
\begin{opts}
  \item 900
  \item 1200
  \item 600
  \item 1500
\end{opts}

\q{Si se fija un precio máximo de 20 en ese mercado, la escasez es:}
\begin{opts}
  \item 20
  \item 30
  \item 40
  \item 50
\end{opts}

\q{Con el precio máximo de 20, la pérdida de bienestar es:}
\begin{opts}
  \item 625
  \item 500
  \item 750
  \item 1250
\end{opts}

\q{Si se aplica un impuesto específico de 5 por unidad, el nuevo precio de equilibrio es:}
\begin{opts}
  \item 30
  \item 31
  \item 32
  \item 33
\end{opts}

\q{Con el impuesto de 5, la cantidad de equilibrio es:}
\begin{opts}
  \item 50
  \item 54
  \item 60
  \item 66
\end{opts}

\q{Con el impuesto de 5, la recaudación total es:}
\begin{opts}
  \item 270
  \item 300
  \item 240
  \item 150
\end{opts}

\q{Con el impuesto de 5, la pérdida de bienestar es:}
\begin{opts}
  \item 5
  \item 10
  \item 30
  \item 15
\end{opts}

\q{En el impuesto anterior, la carga que soportan los consumidores por unidad es:}
\begin{opts}
  \item 2
  \item 4
  \item 3
  \item 5
\end{opts}

\q{En el mismo impuesto, la carga que soportan los productores por unidad es:}
\begin{opts}
  \item 3
  \item 2
  \item 1
  \item 5
\end{opts}

\q{Si se aplica un subsidio de 5 por unidad, el nuevo precio de equilibrio es:}
\begin{opts}
  \item 27
  \item 30
  \item 33
  \item 25
\end{opts}

\q{Con el subsidio de 5, el costo total del subsidio es:}
\begin{opts}
  \item 330
  \item 270
  \item 300
  \item 360
\end{opts}

\q{Con el subsidio de 5, la pérdida de bienestar es:}
\begin{opts}
  \item 5
  \item 10
  \item 15
  \item 30
\end{opts}

\q{En una industria de costo constante, la oferta de largo plazo es:}
\begin{opts}
  \item ascendente.
  \item horizontal.
  \item descendente.
  \item vertical.
\end{opts}

\q{En una industria de costo creciente, la oferta de largo plazo es:}
\begin{opts}
  \item horizontal.
  \item descendente.
  \item ascendente.
  \item vertical.
\end{opts}
